\documentclass[11pt,a4paper]{article}
\usepackage[margin=25mm]{geometry}
\usepackage{amsmath,amssymb,amsthm,hyperref}
\setlength{\emergencystretch}{2em}
\newtheorem{lemma}{Lemma}
\newtheorem{theorem}[lemma]{Theorem}
\newcommand{\gen}[1]{\langle#1\rangle}
\title{Excluding the interior two-bridge equality}
\author{Human Compression project: L044 continuation}
\date{}
\begin{document}\maketitle
Let $G=\gen{\rho_0,\ldots,\rho_{r-1}}\leq A_n$, $n\geq12$,
be transitive with a connected string C-marking and all maximal
omissions intransitive. Assume the high-rank equality configuration
of the accepted nested-bridge theorem, with both crossings interior:
\[
 i=k-1,\quad h=k,\quad j=k+1,\quad
 r=k+l+1,\quad k,l\geq2,\quad n=2r.
\]
The three regions are $A$, $M=\{b,c\}$, $D$, of sizes $2k,2,2l$.
The unique crossing edges are $a\mathrel{-_i}b$ and
$c\mathrel{-_j}d$. Write
\[
 L=G_{<i},\quad s=\rho_i,\quad \tau=\rho_h,
 \quad t=\rho_j,\quad R=G_{>j}.
\]
Here $L$ fixes $M\cup D$ and $R$ fixes $A\cup M$.
We impose the two additional \emph{literal} purity conditions
\[
                         s|_D=1,\qquad t|_A=1.          \tag{1}
\]
These hold in the intended Proposition~5.28 routing when $h$ is
the last shared colour for the $i$-cut and $g=h$ the first shared
colour for the $j$-cut: $t$ already moves $M\cup D$ across its
bridge and therefore cannot move $A$; likewise $s$ already moves
$A\cup M$ and cannot move $D$. We do not infer (1) from the
numerical equalities alone.

The single middle edge is $b\mathrel{-_h}c$. Thus
\[
 \tau|_M=(b\ c),\quad s=(a\ b)u_A,\quad t=(c\ d)v_D,
\]
with $u_A,v_D$ supported inside the indicated outside regions.
The original omission orbits give
\[
              \gen{L,\tau}\text{ transitive on }A,
 \qquad       \gen{\tau,R}\text{ transitive on }D.       \tag{2}
\]
All other generators fix those respective regions, by (1) and
the stated prefix/suffix support conditions.

\begin{lemma}[The exact critical side models]\label{models}
The involution $\tau$ exchanges $k$ pairs in $A$ and $l$ pairs
in $D$. On $A$, either $L$ is transitive, or it has exactly two
orbits of size $k$, each with faithful natural action $S_k$.
On the $k$ pairs of $\tau|_A$, its distinguished $k-1$ generators
form an independent involutory string generating $S_k$. For odd
$k\geq5$ this is the ordinary adjacent-transposition path; for
$k=3$ it is the two-transposition path.
The dual statements hold for $R$ on $D$.
If $k=3$, the transitive alternative for $L$ is impossible.
\end{lemma}
\begin{proof}
The element $\tau$ centralizes $L$. If $L$ is transitive on
$A$, then its nontrivial centralizing involution is fixed-point-free.
The pair action of $L$ is transitive. Every distinguished generator
is essential in that action: if a proper distinguished subset $J$
had the same pair image, $\gen{J,\tau}$ would be transitive on
$A$. A discarded generator is supported on $A$, so adjoining it
to its global omission could not join omission orbits. That omission
would already be transitive on $\Omega$, a contradiction.
Whiston's equality bound therefore makes the $k-1$ pair images
an independent family generating $S_k$.

If $L$ is intransitive, use the accepted faithful commuting-orbit
normalization. Its orbit size is at least $k=i+1$, there are at
least two orbits, and $|A|=2k$. Hence there are exactly two
$k$-point orbits. Whiston's equality theorem makes the faithful
action on each $S_k$, and $\tau$ exchanges the two orbits,
identifying their actions. Again it is fixed-point-free and the
pair action is the same natural $S_k$ action.

In either case the independent involutory pair images form a string.
For odd $k\geq5$, the complete paid L047 maximal-string theorem
and its degree-five certificate make them an ordinary path. For
$k=3$, the two involutions generating $S_3$ are distinct transpositions,
so again give its ordinary path. The subgroup omitting the near
endpoint fixes the pair $\{a,a^\tau\}$ pointwise, since it fixes
$a$ by the unique-crossing lemma and centralizes $\tau$; this identifies
the anchor endpoint of the path. No classification for even $k$ is
needed here, and no projected IP is assumed.

For $k=3$ in the transitive alternative, identify the paired action
with a subgroup of $C_2^3\rtimes S_3$. Both distinguished generators
are even on six points, so their binary flip vectors have even sum.
Their product has pair image a $3$-cycle. Its cube is a constant
flip vector whose coordinate is the sum of the product's flip vector,
and that sum is zero. Thus their product has order exactly three.
Two involutions with product of order three generate $S_3$ of
order six. A transitive action of that group on six points is
regular, whereas its nonidentity far generator fixes $a$.
This contradiction leaves precisely the two natural three-point
orbits. Every argument reverses for the right side.
\end{proof}

\begin{lemma}[No odd critical side]\label{odd}
Under the standing hypotheses and (1), $k$ is even. By reversal,
$l$ is even as well.
\end{lemma}
\begin{proof}
First suppose $k\geq5$ is odd. Let $H=G_{<i-1}$, the omission
of the near endpoint of $L$. By Lemma~\ref{models}, $H$ fixes
$a$ and $a'=a^\tau$, and induces the full $S_{k-1}$ on the
remaining $k-1$ pairs of $A$. Those two anchors are its only
fixed points in $A$. Since $s=(a\ b)u_A$ centralizes $H$ and
$u_A$ preserves $A\setminus\{a\}$, it fixes $a'$. Its action
on the other $2(k-1)$ points centralizes $H$. The complete
accepted paired-row centralizer lemma, with even $k-1\geq4$,
says that this involution is even. Hence $u_A$ is even and
$s$ is odd, contradicting $G\leq A_n$.

It remains to exclude $k=3$. By Lemma~\ref{models}, after
identifying the two natural rows through $\tau$, write
\[
 A=\{1,2,3,4,5,6\},\quad a=3,\quad a'=6,
\]
\[
 x=\rho_0=(1\ 2)(4\ 5),\quad
 y=\rho_1=(2\ 3)(5\ 6),\quad
 \tau|_A=(1\ 4)(2\ 5)(3\ 6).
\]
Here $i=2$, so $s=(3\ b)u_A$ commutes with $x$. The point
$6$ is fixed by $u_A$: it is the sole $x$-fixed point in
$A\setminus\{3\}$. On $\{1,2,4,5\}$ an involution
centralizing $(1\ 2)(4\ 5)$ either preserves both pairs or
exchanges them. In the latter case it has two transpositions
and is even. In the former case its odd possibilities are
exactly one pair flip. Evenness of $s$ therefore forces
\[
                    u_A=(1\ 2)\quad\hbox{or}\quad(4\ 5).       \tag{3}
\]

Now omit $x$ from the \emph{global} group. The remaining
$\tau,t,R$ are transitive on $M\cup D$: equation (2) makes
$\gen{\tau,R}$ transitive on $D$, the edge $c\mathrel{-_t}d$
joins $c$, and $b\mathrel{-_\tau}c$ joins $b$.
The $s$-edge $3b$ and the $\tau$-edge $36$ then add $3,6$
to this omission orbit. The $y$-edges add $2,5$. In the first
case of (3) the other $s$-edge adds $1$, and $\tau$ adds $4$;
in the second case it adds $4$, and $\tau$ adds $1$.
Thus the omission is transitive on all $\Omega$, contradicting
our standing hypothesis. This excludes three, and every odd
$k\geq2$ has now been excluded. Reversing the entire argument
proves the assertion for $l$.
\end{proof}

\begin{theorem}[Interior touching bands are impossible]
The equality $g=h$ in the nested-bridge high-rank theorem cannot
occur with both unique crossings interior and with the literal
purity conditions (1). In particular this equality is excluded
when $h$ and $g$ are the last and first shared colours in the
oriented Proposition~5.28 configuration.
\end{theorem}
\begin{proof}
Lemma~\ref{models} and $M=\{b,c\}$ show that $\tau$ is a
product of exactly $k+l+1$ transpositions. Lemma~\ref{odd}
makes both $k,l$ even. Thus $\tau$ is odd, the final contradiction.
\end{proof}

\paragraph{Attribution, dependencies and boundary.}
The source is Cameron--Fernandes--Leemans--Mixer, \emph{Highest
rank of a polytope for $A_n$}, arXiv:1605.09173v1,
Proposition~5.28, physical page~37. The new proof is the critical
side parity and literal omission argument above. It reuses the
complete accepted nested-bridge count, commuting-orbit and pair
image arguments, the whole L047 maximal-string theory and degree-five
certificate, Whiston's bounds, and the paired-row centralizer lemma.
No sample-only claim or new finite computation is a premise.
The degree-six norm calculation is proved here directly.

Endpoint touching configurations, the separate $g>h$ middle-orbit
equality and the remaining full-source routing and rigidity are
not closed by this note. There is no new formal declaration,
whole-original recovery or question/priority claim.
\end{document}
